%%======================================================================
%%  Chapitre 4 : Réalisation du projet — Rapport 2 (Lot B)
%%======================================================================
\chapter{Réalisation du projet — Rapport 2 (Lot B)}

Ce chapitre présente la \textbf{Réalisation Lot B} du même projet \textit{EduPlatform}. Le périmètre couvre principalement le \textbf{frontend}, l'\textbf{expérience utilisateur} et les \textbf{parcours Enseignant/Étudiant}.

%%======================================================================
\section{Périmètre du Lot B}

Le Lot B regroupe les tâches suivantes :
\begin{itemize}[noitemsep]
    \item réalisation des interfaces web React/TypeScript ;
    \item implémentation des espaces \textbf{enseignant} et \textbf{étudiant} ;
    \item amélioration UX (responsive, mode sombre, navigation rapide) ;
    \item intégration côté client des endpoints backend et validation des parcours utilisateurs.
\end{itemize}

%%======================================================================
\section{Interfaces communes et authentification}

La couche frontend a été conçue pour proposer une expérience cohérente :
\begin{itemize}[noitemsep]
    \item \textbf{Landing Page} informative avec sections dynamiques ;
    \item \textbf{Login multi-rôle} (Admin, Enseignant, Étudiant) ;
    \item routage protégé et redirection vers les dashboards par rôle ;
    \item persistance de session et préférences utilisateur.
\end{itemize}

%%======================================================================
\section{Réalisation de l'espace Enseignant}

Les fonctionnalités réalisées pour l'enseignant incluent :
\begin{itemize}[noitemsep]
    \item dashboard de suivi pédagogique ;
    \item gestion des cours et des ressources ;
    \item planification des évaluations ;
    \item consultation de l'emploi du temps ;
    \item traitement des réclamations de notes ;
    \item émission de notifications soumises à approbation.
\end{itemize}

%%======================================================================
\section{Réalisation de l'espace Étudiant}

Les fonctionnalités majeures réalisées côté étudiant :
\begin{itemize}[noitemsep]
    \item consultation des cours et téléchargement des ressources ;
    \item gestion des devoirs et soumissions ;
    \item consultation des notes et moyenne ;
    \item suivi des absences et dépôt de justifications ;
    \item demandes administratives et demandes de stage ;
    \item réclamations et notifications.
\end{itemize}

%%======================================================================
\section{Qualité UX et validation fonctionnelle}

Le Lot B a également couvert :
\begin{itemize}[noitemsep]
    \item responsive design (mobile/tablette/desktop) ;
    \item mode sombre avec persistance ;
    \item navigation rapide (palette de commandes) ;
    \item vérification de bout en bout des scénarios utilisateur principaux.
\end{itemize}

%%======================================================================
\section{Répartition des tâches — Rapport 2}

\begin{table}[H]
\centering
\caption{Division des tâches de réalisation — Lot B}
\label{tab:taches_realisation_lot_b}
\begin{tabularx}{\textwidth}{lX}
\toprule
\textbf{Axe} & \textbf{Tâches prises en charge} \\
\midrule
Frontend & Développement des composants React, pages par rôle et intégration API côté client. \\
Parcours Enseignant & Cours, évaluations, emploi du temps, réclamations de notes, notifications. \\
Parcours Étudiant & Devoirs/soumissions, notes, absences/justifications, demandes et réclamations. \\
UX/UI & Responsive design, mode sombre, ergonomie de navigation et feedback utilisateur. \\
Validation fonctionnelle & Vérification des parcours métier côté interface et cohérence des interactions client-serveur. \\
\bottomrule
\end{tabularx}
\end{table}

\newpage
